% chapters/07-testing.tex -- part B. 시험 방법과 결과.
% Macros come from part A's preamble.tex. \code is verbatim-like: not in \caption, \section, node text.
\chapter{시험 방법과 결과}\label{ch:testing}

% src: deploy/traffic/review/02-blocking.md §시험 결과; traffic.design.md §10
\paragraph{결과.}
2026/10/13에 시험 네 가지를 돌렸다.
캐시 \cfs{HIT} 경로(L1)는 640건이 모두 \cfs{HIT}였고 VM에 닿은 시험 요청은 0건이었다.
오리진(origin) \cfs{MISS} 경로는 웹이 초당 2번까지 버티고 4번에서 시간 초과가 생겼다.
API는 초당 8번에서도 느려지지 않았다.
Cloudflare 요청 수 제한은 한도 100과 50에서 모두 429를 냈다.
\Cref{tab:07-results-summary}가 한눈에 보인다.

\begin{table}[H]
  \caption{시험 결과 요약 (2026/10/13). 자세한 숫자는 이 장의 각 절에 있다.}\label{tab:07-results-summary}
  % src: deploy/traffic/review/02-blocking.md §시험 결과 L1, L2, L-a, L2-api, L-a 한도 50
  \small
  \begin{tabularx}{\linewidth}{L{2.6cm} Y L{2.2cm}}
    \toprule
    시험 & 한 줄 결과 & 판정 \\
    \midrule
    L1 (run \texttt{38465971720}) & 640건 모두 \cfs{HIT}. VM까지 온 시험 요청 0건 & 통과 \\
    \midrule
    L2 web (run \texttt{38466353765}) & 4 rps 단계에서 10건 시간 초과 (3.1\,\%). 5xx 0 & 중단 \\
    \midrule
    L2-api (run \texttt{38468411633}) & 450건 모두 \cfs{MISS}, 실패 0, 5xx 0, p95 626 ms & 통과 \\
    \midrule
    Cloudflare 한도 100 & 250번 중 200 232, 429 18 & 발동 \\
    \midrule
    Cloudflare 한도 50 & 150번 중 200 53, 429 97 & 발동 \\
    \bottomrule
  \end{tabularx}
\end{table}

%% ---------------------------------------------------------------------------
\section{시험 철학: 고의 파손}\label{sec:07-philosophy}

\begin{inbrief}
  \item[Why] 단위 시험을 쓰지 않는 저장소에서, 규칙이 정말 지켜지는지 증명할 방법이 필요하다.
  \item[When] 오너 지시(2026/09/04)로 시험이 중지된 동안이다.
  \item[Where] 실제 운영 시스템 하나뿐이다.
  \item[How] 일부러 고장을 일으키고, 관찰하고, 정리한다.
\end{inbrief}

% src: .claude/rules/core.md §Prohibitions (Tests are SUSPENDED, owner directive 2026-09-04)
이 저장소에서 \code{*.test.*}와 \code{*.spec.*} 파일은 쓰지 않는다.
시험 실행기(test runner)도 돌리지 않는다.
대신 실제 시스템에 \term{고의 파손}\eng{deliberate breakage}을 한다.
고장을 일으키고, 결과를 보고, 정리한다.

\begin{figure}[H]
  \centering
  \begin{tikzpicture}[font=\small]
    \node[flowbox, minimum width=3.6cm] (a) at (0,0) {1. 고장을\\일으킨다};
    \node[flowbox, minimum width=3.6cm] (b) at (5.2,0) {2. 결과를\\관찰한다};
    \node[flowbox, minimum width=3.6cm] (c) at (10.4,0) {3. 정리한다};
    \draw[flowarrow] (a) -- (b);
    \draw[flowarrow] (b) -- (c);
  \end{tikzpicture}
  \caption{고의 파손의 세 걸음. 규칙 원문은 ``trigger the failure, observe it, clean up''이다.
  (src: .claude/rules/core.md §Prohibitions)}\label{fig:07-breakage-steps}
\end{figure}

\Cref{tab:07-test-kinds}가 이 장의 시험 세 종류를 나눈다.

\begin{table}[H]
  \caption{시험 세 종류. 무엇을 알려고 쓰는가.}\label{tab:07-test-kinds}
  % src: deploy/traffic/README.md §Monitoring commands; deploy/traffic/load/README.md §1; traffic.design.md §10.1
  \small
  \begin{tabularx}{\linewidth}{L{3.6cm} Y L{3.4cm}}
    \toprule
    시험 & 무엇을 알려고 쓰나 & 어디서 도나 \\
    \midrule
    \code{traffic-check.sh} & 캐시 규칙이 아직 맞는가. 개인 응답이 캐시되지 않는가 & Mac, GitHub runner \\
    \midrule
    고의 파손 & 한 가지 규칙이 정말 막는가 (403, 429 등) & Mac의 \code{curl} \\
    \midrule
    k6 부하 시험 & 엣지가 몰리는 요청을 흡수하는가. 오리진은 \cfs{MISS}를 몇 rps까지 버티는가. 제한이 발동하는가 & GitHub runner만 \\
    \bottomrule
  \end{tabularx}
\end{table}

%% ---------------------------------------------------------------------------
\section{traffic-check.sh: 14개 검사}\label{sec:07-trafficcheck}

\begin{inbrief}
  \item[Why] 캐시 규칙이 나중에 깨져도 사람이 매번 확인하지 않아도 알게 한다.
  \item[When] 손으로 필요할 때, 부하 시험 시작 전, 매주 월요일에 돈다.
  \item[Where] \code{_scripts/cloudflare/traffic-check.sh}. 읽기 전용, GET 약 20번이다.
  \item[How] 검사마다 기대 패턴(정규식)과 실제 값을 비교해 PASS나 FAIL을 찍는다. FAIL이 있으면 종료 코드 1이다.
\end{inbrief}

% src: _scripts/cloudflare/traffic-check.sh; deploy/traffic/README.md §Monitoring commands §2 c, §3 c
\paragraph{결과.}
2026/10/13에 14개 모두 PASS였다 (0 FAIL).
\Cref{tab:07-checks}는 검사 14개의 기대값과 README에 적힌 측정값이다.

\begin{table}[H]
  \caption{\texttt{traffic-check.sh}의 검사 14개. 기대값은 스크립트의 정규식이다.}\label{tab:07-checks}
  % src: _scripts/cloudflare/traffic-check.sh (기대값); deploy/traffic/README.md §2 c Measured, §3 c Measured, §Monitoring 1 (측정값); 01.status.md 2026-10-13 (14/14)
  \footnotesize
  \begin{tabularx}{\linewidth}{r L{3.7cm} L{2.9cm} L{2.9cm} Y}
    \toprule
    \# & 검사 & 기대값 & 측정값 (README) & 결과 \\
    \midrule
    1 & \texttt{GET /} (HIT, 4번 안에) & \cfs{HIT} & \cfs{MISS} 다음 \cfs{HIT} & PASS \\
    \midrule
    2 & \texttt{GET /blog} & \cfs{HIT} & \cfs{HIT} & PASS \\
    \midrule
    3 & \texttt{GET /robots.txt} & \cfs{HIT} & \cfs{HIT} & PASS \\
    \midrule
    4 & \texttt{GET /sitemap.xml} & \cfs{HIT} & \cfs{HIT} & PASS \\
    \midrule
    5 & \code{__session} 쿠키 & \cfs{BYPASS} 또는 \cfs{DYNAMIC} & \cfs{DYNAMIC} & PASS \\
    \midrule
    6 & \code{/me} (보호 경로) & \cfs{BYPASS} 또는 \cfs{DYNAMIC} & 307, \cfs{DYNAMIC} & PASS \\
    \midrule
    7 & \code{rsc} 헤더, \code{_rsc=} 없음 & \cfs{BYPASS} 또는 \cfs{DYNAMIC} & \cfs{DYNAMIC} & PASS \\
    \midrule
    8 & \code{?_rsc=} content-type & \code{text/x-component} & \code{text/x-component} & PASS \\
    \midrule
    9 & HTML 브라우저 cache-control & \texttt{private, no-cache, no-store} 로 시작 & 기록 없음 & PASS \\
    \midrule
    10 & HTML \code{set-cookie} 개수 & 0 & 0 & PASS \\
    \midrule
    11 & \texttt{GET posts?limit=1} (HIT, 4번 안에) & \cfs{HIT} & \cfs{MISS} 다음 \cfs{HIT} & PASS \\
    \midrule
    12 & \code{Authorization} 헤더 & \cfs{BYPASS} 또는 \cfs{DYNAMIC} & \cfs{DYNAMIC} & PASS \\
    \midrule
    13 & \code{/ping} (Cache-Control 없음) & \cfs{BYPASS}, \cfs{DYNAMIC}, \cfs{MISS} 중 하나 & \cfs{BYPASS} & PASS \\
    \midrule
    14 & CORS \code{my.} origin 되돌림 & 요청한 origin과 같음 & \code{my.}와 \code{local.} 모두 자기 \texttt{access-control-}\allowbreak\texttt{allow-origin} & PASS \\
    \bottomrule
  \end{tabularx}
\end{table}

\begin{description}
  \item[읽는 법] 줄마다 PASS나 FAIL이다. 마지막 줄은 \texttt{== N FAIL}이다.
  \item[FAIL의 뜻] 일반 \texttt{GET} 검사가 \cfs{DYNAMIC}으로 FAIL이면 캐시 규칙이 꺼졌거나 식이 바뀐 것이다.
    \code{__session}, \code{/me}, \code{Authorization} 검사가 \cfs{HIT}로 FAIL이면 개인 응답이 캐시된다는 뜻이다. 그때는 ``web cache'' 또는 ``api cache''를 바로 끈다.
    HTML cache-control 검사가 FAIL이면 zone의 Browser Cache TTL이 ``Respect Existing Headers''가 아니다.
\end{description}

%% ---------------------------------------------------------------------------
\section{고의 파손 예}\label{sec:07-breakage-examples}

\begin{inbrief}
  \item[Why] 규칙이 막는다는 증거는 ``막힌 것을 본 것''뿐이다.
  \item[When] 변경을 배포한 직후다.
  \item[Where] 운영 시스템에 읽기 위주로, 한 요청씩이다.
  \item[How] 나쁜 요청 하나를 보내고 기대한 거절 코드가 나오는지 본다.
\end{inbrief}

% src: deploy/traffic/README.md §Monitoring commands (Verified 2026-10-13), §4 b Check; deploy/traffic/review/03-status-codes.md §G §L; deploy/traffic/review/04-observability.md §F 확인; deploy/traffic/review/02-blocking.md §L-a
\Cref{tab:07-breakage}에 실제로 한 고의 파손과, 일부러 하지 않은 것 하나를 적었다.

\begin{table}[H]
  \caption{고의 파손 예. 마지막 줄은 일부러 하지 않은 시험이다.}\label{tab:07-breakage}
  % src: 각 행의 출처는 본문 위 % src 줄
  \footnotesize
  \begin{tabularx}{\linewidth}{L{3.2cm} L{3.5cm} L{2.6cm} Y}
    \toprule
    규칙 & 일부러 한 고장 & 기대 & 관찰한 결과 \\
    \midrule
    캐시 검사 스크립트 (F) & \code{W=https://blog-api.prettylog.com}, 웹 자리에 API 호스트 (일부러 틀린 주소) & 웹 검사 FAIL & 웹 검사가 FAIL을 찍었다 \\
    \midrule
    CORS 거절 (G) & 허용 안 된 Origin의 OPTIONS와 GET (\code{https://evil.example}) & 403 & 403. \code{/users/me}는 \texttt{origin not allowed}. ERROR 0건 \\
    \midrule
    error-logs 키 (L) & 키 없이 POST, 틀린 키로 POST & 403 & 둘 다 403 (\texttt{forbidden}), INSERT 없음 \\
    \midrule
    JSON 404 (제안 2) & 없는 경로 \code{/v1/api/nope} & 404 JSON & 404 JSON \\
    \midrule
    Cloudflare 요청 수 제한 (L-a) & \code{/}를 250번, 6초, 동시 25 & 429 & 200 232, 429 18 \\
    \midrule
    메모리 알림 (E) & \textbf{하지 않음} & (없음) & 하지 않았다. VM이 1 GB라서 일부러 채우면 운영 사이트가 위험하다 \\
    \bottomrule
  \end{tabularx}
\end{table}

CORS 시험에는 뒷이야기가 있다.
주간 issue \#5의 ERROR 48줄은 24건이다.
그중 23건은 로컬 설정이고 1건이 고의 파손 시험이었다.
모두 우리가 만든 요청이다.
메모리 알림은 고의 파손의 한계를 보인다.
운영 사이트를 위험하게 하는 고의 파손은 하지 않는다.
대신 정책이 있고 켜져 있다는 사실만 확인한다.

%% ---------------------------------------------------------------------------
\section{요청 수 제한 설계 (H)}\label{sec:07-ratelimit-design}

\begin{inbrief}
  \item[Why] 캐시는 익명 GET만 막는다. 캐시를 지난 요청이 몰리면 1 GB VM이 버티지 못한다.
  \item[When] 2026/10/13에 설계했다. L-a는 켰고, L-b는 \texttt{log} 모드다.
  \item[Where] Cloudflare 엣지(L-a)와 Express(L-b)다.
  \item[How] 두 겹으로 센다. 시험은 이 겹들이 정말 발동하는지 본다.
\end{inbrief}

% src: deploy/traffic/review/02-blocking.md §H; decisions.plan.md #62 #67 #68
\paragraph{결과.}
Cloudflare 한도는 IP당 100에서 50 / 10 s로 낮췄다 (\decision{68}).
Express 한도는 읽기 120 / 60 s, 쓰기 30 / 60 s이고 처음 14일은 세기만 한다 (\decision{67}).

\begin{table}[H]
  \caption{요청 수 제한 세 층.}\label{tab:07-layers}
  % src: deploy/traffic/review/02-blocking.md §H 무엇을 바꾸나 표; decisions.plan.md #62 #68
  \small
  \begin{tabularx}{\linewidth}{L{2.4cm} L{3cm} L{3.8cm} Y}
    \toprule
    층 & 대상 & 한도 & 동작과 끄는 법 \\
    \midrule
    L-a Cloudflare & \code{my.} & IP당 50 / 10 s (처음 가안 100, \decision{68}) & Block 10 s. 대시보드 토글 \\
    \midrule
    L-b Express & \code{blog-api} & 읽기 120 / 60 s, 쓰기 30 / 60 s & 429. 처음 14일은 \texttt{log} 모드(세기만 함). \code{RATE_LIMIT_MODE=off} 후 ecosystem restart \\
    \midrule
    L-c & HIT 경로, SSR & 없음 & 없음 \\
    \bottomrule
  \end{tabularx}
\end{table}

\begin{itemize}[nosep]
  \item \code{blog-api.}를 엣지에서 세지 않는다. Next SSR이 API를 부르면 모두 VM IP 하나로 보인다. 엣지에서 IP로 세면 사이트가 자기 SSR을 막는다.
  \item SSR 요청은 \texttt{x-origin-key}가 맞으면 세지 않는다 (\decision{64}).
  \item Challenge 대신 Block이다. RSC 요청과 \code{fetch}는 challenge를 풀 수 없다. 앱 안 이동이 깨진다.
  \item 카운터는 프로세스 메모리(Map)에 둔다. DB나 Redis는 요청마다 쓰기가 늘고 VM 한 대에 이득이 없다 (\decision{62}).
\end{itemize}

%% ---------------------------------------------------------------------------
\section{k6 부하 시험: 방법}\label{sec:07-k6}

\begin{inbrief}
  \item[Why] 세 가지를 알려고 한다. 엣지가 burst를 흡수하나, 오리진은 \cfs{MISS}를 몇 rps까지 버티나, 제한이 발동하나.
  \item[When] 오너가 세트마다 승인한 뒤에, 가장 한가한 시간대에 돈다.
  \item[Where] GitHub Actions의 \code{load-test.yml}에서만 돈다. Mac에서는 돌리지 않는다.
  \item[How] k6가 정해진 속도로 공개 URL에 GET을 보내고, 중단 조건에 걸리면 스스로 멈춘다.
\end{inbrief}

% src: deploy/traffic/load/README.md §1; traffic.design.md §10.1 §10.2
\subsection{세트별 목적}\label{sec:07-k6-sets}

\begin{table}[H]
  \caption{부하 시험 세트. 무엇을 알려고 쓰나.}\label{tab:07-sets}
  % src: deploy/traffic/load/README.md §1 표; traffic.design.md §10.2 표 (L3)
  \footnotesize
  \begin{tabularx}{\linewidth}{L{1.6cm} L{2.8cm} L{3.8cm} L{3cm} Y}
    \toprule
    세트 & 답하는 질문 & 대상 & 속도 $\times$ 시간 & 총 요청 \\
    \midrule
    L1 (\code{l1.js}) & 엣지가 burst를 흡수하는가 & \code{my.prettylog.com/blog/@taeklim/rust-config}, \code{blog-api.prettylog.com/v1/api/posts?limit=1} (둘 다 \cfs{HIT}) & 각 8 rps $\times$ 40 s & 320 + 320 = 640 \\
    \midrule
    L2 (\code{l2.js}) web & 오리진은 \cfs{MISS}를 몇 rps까지 버티는가 & \code{my.prettylog.com/?lt=<run>-<n>} (매번 \cfs{MISS}) & 0.5 $\rightarrow$ 1 $\rightarrow$ 2 $\rightarrow$ 4 rps, 단계당 60 s & 450 \\
    \midrule
    L2-api (\code{l2.js}의 api 절반) & 같은 질문, API 쪽 & \code{blog-api.prettylog.com/v1/api/posts?limit=1&lt=<run>-<n>} (매번 \cfs{MISS}) & 1 $\rightarrow$ 2 $\rightarrow$ 4 $\rightarrow$ 8 rps, 단계당 30 s & 450 \\
    \midrule
    L3-cf (계획) & Cloudflare 제한이 발동하는가 & \code{my.prettylog.com/?lt=<run>-<n>} & 한도를 잠시 20 / 10 s로 낮춘 뒤 3 rps $\times$ 20 s & 60 \\
    \midrule
    L3-api (계획) & Express 제한이 발동하는가 & \code{blog-api.prettylog.com/} (DB 없음, 로그 남음) & 2.5 rps $\times$ 120 s & 300 \\
    \bottomrule
  \end{tabularx}
\end{table}

\begin{itemize}[nosep]
  \item L2는 web을 먼저, api를 그 뒤에 돌린다. 두 오리진에 동시에 부하를 주지 않는다.
  \item \code{lt=} 값이 매번 달라서 캐시 키(cache key)가 \cfs{MISS}가 된다. 이 항목은 하루 뒤 만료되고 purge하지 않는다.
  \item 모든 요청의 User-Agent는 \code{prettylog-load-test/L1} 또는 \code{prettylog-load-test/L2}다. Cloud Logging에서 걸러 낼 수 있다.
  \item L3를 창(window) 길이의 2배로 보내는 이유: 고정 창은 시계 경계에서 끊긴다. 한 창 길이만 보내면 두 창으로 나뉘어 한도를 안 넘을 수 있다.
\end{itemize}

\Cref{fig:07-l2-ramp}는 L2 web 단계를 그림으로 보인다.
각 60초 단계 끝에서 중단 조건을 평가한다.
통과한 마지막 단계의 rps가 C 값이다.

\begin{figure}[H]
  \centering
  \begin{tikzpicture}[font=\footnotesize]
    \draw[->] (0,0) -- (9.4,0) node[right] {시간};
    \draw[->] (0,0) -- (0,4.2) node[above] {rps};
    \foreach \y/\l in {0.7/0.5, 1.4/1, 2.8/2, 3.8/4}
      \draw (0.08,\y) -- (-0.08,\y) node[left] {\l};
    \draw[very thick, vmcol] (0,0.7) -- (2.2,0.7) -- (2.2,1.4) -- (4.4,1.4) -- (4.4,2.8) -- (6.6,2.8) -- (6.6,3.8) -- (8.8,3.8);
    \foreach \x/\l in {1.1/{단계 1\\60 s}, 3.3/{단계 2\\60 s}, 5.5/{단계 3\\60 s}, 7.7/{단계 4\\60 s}}
      \node[align=center, below] at (\x,-0.1) {\l};
    \node[badcol, align=center, font=\footnotesize] at (7.7,4.4) {실제 run은\\여기서 중단};
    \node[goodcol, align=center, font=\footnotesize] at (5.5,3.35) {C\textsubscript{web} 약 2 rps};
  \end{tikzpicture}
  \caption{L2 web 단계 계획. 실제 run은 4 rps 단계에서 10건 시간 초과로 중단했다.
  (src: traffic.design.md §10.3, deploy/traffic/review/02-blocking.md §L2)}\label{fig:07-l2-ramp}
\end{figure}

\subsection{중단 조건과 통과 기준}\label{sec:07-k6-criteria}

중단 조건은 k6 \texttt{thresholds}의 \texttt{abortOnFail}이다.
시작 10초 뒤부터 평가한다 (\texttt{delayAbortEval}).
기준은 run 전에 정했고, run 뒤에 바꾸지 않는다 (사전 등록).

\begin{table}[H]
  \caption{중단 조건.}\label{tab:07-stop}
  % src: deploy/traffic/load/README.md §2; traffic.design.md §10.4 item 5
  \small
  \begin{tabularx}{\linewidth}{L{2cm} Y}
    \toprule
    세트 & 중단 조건 \\
    \midrule
    L1 & 200이 아닌 응답이 1\,\%를 넘는다 \\
    \midrule
    L2 & 5xx가 1건 이상, 실패율이 2\,\%를 넘는다, p95가 3 s를 넘는다. 하나라도 걸리면 중단 \\
    \midrule
    L3 & 없음 (429가 기대값이다) \\
    \bottomrule
  \end{tabularx}
\end{table}

\begin{table}[H]
  \caption{통과 기준 (사전 등록, 설계 \S10.5).}\label{tab:07-pass}
  % src: traffic.design.md §10.5 표
  \small
  \begin{tabularx}{\linewidth}{L{2.8cm} Y}
    \toprule
    세트 & 통과 \\
    \midrule
    L1 & 응답 200 $\geq$ 99\,\%, \cfs{HIT} $\geq$ 99\,\%, run 시간대 blog-web의 그 path 줄 $\leq$ 10, PM2 재시작 횟수 불변 \\
    \midrule
    L2 & 단계별 판정. 재시작 횟수 불변, 메모리 최댓값 $<$ 80\,\%, blog-api 5xx 0 (L2-api) \\
    \midrule
    L3-cf & 60건 중 10건 이상 차단, blog-web 그 path 줄 $<$ 60, 한도 100으로 복구 확인 \\
    \midrule
    L3-api \texttt{log} & 300건 모두 200, \texttt{rate\_limit} warn 줄 $\geq$ 30 \\
    \midrule
    L3-api \texttt{enforce} & 429 + \texttt{Retry-After} $\geq$ 30건. 절반은 \texttt{x-origin-key: wrong}을 붙여도 429. \texttt{error\_log} 증가 0 \\
    \midrule
    SSR 제외 & \texttt{log} 모드 14일 동안 \texttt{req.ssr=true} 줄이 있고, 그중 \texttt{rate\_limit} warn은 0 \\
    \bottomrule
  \end{tabularx}
\end{table}

\begin{itemize}[nosep]
  \item L2의 첫 run은 판정하지 않는다. C\textsubscript{web}, C\textsubscript{api}의 첫 값만 기록한다.
  \item 메모리 80\,\%인 이유: 경보는 90\,\%가 5분 이어질 때다. 단계는 60초라서 경보가 울리지 않을 수 있다. 경보는 통과 신호가 아니다.
  \item k6는 VM 메모리를 볼 수 없다. 그래서 메모리와 재시작 횟수는 run 뒤 점검이 판정한다.
\end{itemize}

\subsection{workflow가 도는 길}\label{sec:07-k6-workflow}

% src: .github/workflows/load-test.yml; deploy/traffic/load/README.md §3 §4 §5; traffic.design.md §10.4
\begin{figure}[H]
  \centering
  \begin{tikzpicture}[font=\footnotesize]
    \node[flowbox, minimum width=3cm] (cmd) at (0,0) {main session\\\texttt{gh workflow run}\\\texttt{load-test.yml}\\\texttt{-f set=...}};
    \node[flowbox, minimum width=3cm] (pre) at (4.2,0) {사전 점검\\\texttt{traffic-check.sh}\\0 FAIL이어야 시작\\(L1은 2번째 GET \cfs{HIT})};
    \node[flownew, minimum width=3cm] (k6) at (8.4,0) {k6 2.3.0 실행\\중단 조건에 걸리면\\스스로 멈춤};
    \node[cfbox, minimum width=2.6cm] (cf) at (12.4,0) {Cloudflare 엣지\\\cfs{HIT}면 응답};
    \node[vmbox, minimum width=2.6cm] (vm) at (12.4,-2.6) {VM\\\cfs{MISS}만 도착\\blog-web, blog-api};
    \node[dbbox, minimum width=2.6cm] (db) at (12.4,-5.0) {Supabase\\API \cfs{MISS} 때만};
    \node[flowbox, minimum width=3cm] (out) at (8.4,-2.6) {\texttt{summary.md}\\$\rightarrow$ Job summary\\\texttt{summary.json}\\$\rightarrow$ artifact};
    \begin{scope}[on background layer]
      \node[zone, draw=mutecol, fit=(pre)(k6), label={[font=\scriptsize,mutecol]above:GitHub runner (ubuntu-latest)}] {};
    \end{scope}
    \draw[flowarrow] (cmd) -- (pre);
    \draw[flowarrow] (pre) -- (k6);
    \draw[flowarrow] (k6) -- (cf) node[lbl, midway, above] {GET};
    \draw[flowarrow] (cf) -- (vm);
    \draw[flowarrow] (vm) -- (db);
    \draw[flowarrow] (k6) -- (out);
  \end{tikzpicture}
  \caption{\texttt{load-test.yml}이 도는 길. runner는 공개 URL만 부르고 비밀 값을 쓰지 않는다.
  (src: .github/workflows/load-test.yml, traffic.design.md §10.4)}\label{fig:07-k6-workflow}
\end{figure}

\begin{table}[H]
  \caption{workflow와 k6의 상한. 시험이 서버를 해치지 않게 하는 장치다.}\label{tab:07-caps}
  % src: .github/workflows/load-test.yml; deploy/traffic/load/common.js; deploy/traffic/load/l2.js; traffic.design.md §10.4 item 4
  \small
  \begin{tabularx}{\linewidth}{L{4.6cm} Y}
    \toprule
    상한 & 값 \\
    \midrule
    실행 시간 & \texttt{timeout-minutes: 8}. L2는 약 6.5분 필요 \\
    \midrule
    동시 실행 & \texttt{concurrency: load-test}, 1개 (\texttt{cancel-in-progress: false}) \\
    \midrule
    입력 & 세트 이름 선택지(L1, L2, L2-api)뿐. 자유 URL 입력 없음 \\
    \midrule
    비밀 값 & 필요 없음. 공개 URL만 호출 \\
    \midrule
    k6 VU 상한 & \texttt{maxVUs} 20, 미리 할당 10 \\
    \midrule
    요청당 \texttt{timeout} / 리다이렉트 & 10 s / 0 \\
    \midrule
    총 요청 상한 & 스크립트의 \texttt{overCap}이 L1 대상당 320, L2 대상당 450으로 막음 \\
    \midrule
    성공으로 치는 응답 & 200만 (리다이렉트나 4xx는 실패로 센다) \\
    \midrule
    k6 버전 & 2.3.0 \\
    \bottomrule
  \end{tabularx}
\end{table}

\subsection{실행 순서}\label{sec:07-k6-run}

\begin{enumerate}[nosep]
  \item 시작 전: L-b가 \texttt{off}나 \texttt{log}인지 본다. L2-api 8 rps는 \texttt{enforce}면 막힌다.
  \item 시작 전: PM2 재시작 횟수를 적어 둔다. 가장 한가한 시간대를 고른다.
  \item 실행: \texttt{gh workflow run load-test.yml -f set=L1}. 세트는 L1, L2, L2-api 중 하나다.
  \item 관찰: \texttt{gh run watch <id>}. 이상하면 \texttt{gh run cancel <id>}로 바로 끊는다.
  \item 끝난 뒤: 재시작 횟수가 같은지, 메모리 최댓값(\texttt{./\_scripts/gcp/vm-metrics.sh 1}), blog-web 로그의 시험 줄 수, blog-api 로그의 5xx 수를 본다.
\end{enumerate}

원본 숫자는 artifact \code{load-test-<set>-<id>}의 \texttt{summary.json}이다.
이 장의 숫자는 그 요약과 README에 기록된 값을 옮겼다.

%% ---------------------------------------------------------------------------
\section{결과 1: L1, 캐시 \texorpdfstring{\cfs{HIT}}{HIT} 경로}\label{sec:07-result-l1}

\begin{inbrief}
  \item[Why] 익명 GET 폭주를 엣지가 받아 주는지 본다.
  \item[When] run \texttt{38465971720}, 2026/10/13.
  \item[Where] \code{my.}의 글 페이지와 \code{blog-api.}의 \code{posts?limit=1}이다.
  \item[How] 각 8 rps $\times$ 40 s. 640건을 보낸다.
\end{inbrief}

% src: deploy/traffic/review/02-blocking.md §L1: 캐시 HIT 경로 (run 38465971720, 통과)
\paragraph{결과.}
통과했다.
web 320건, api 320건이 모두 \cfs{HIT}였다.
실패 0건이고 VM에 닿은 시험 요청은 0건이었다.

\begin{table}[H]
  \caption{L1 결과 (run \texttt{38465971720}).}\label{tab:07-l1}
  % src: deploy/traffic/review/02-blocking.md §L1 표
  \small
  \begin{tabularx}{\linewidth}{L{3.6cm} Y}
    \toprule
    항목 & 결과 \\
    \midrule
    요청 & 640 (web 320, api 320), 8 rps $\times$ 40 s 씩 \\
    \midrule
    \texttt{cf-cache-status} & web \cfs{HIT} 320, api \cfs{HIT} 320 \\
    \midrule
    실패 & 0 \\
    \midrule
    p95 & 전체 82 ms (web 88, api 73) \\
    \midrule
    VM까지 온 시험 요청 & 0 (blog-web 로그에 \texttt{prettylog-load-test} 없음) \\
    \midrule
    PM2 재시작 & 전후 같음 (api 5, web 3) \\
    \midrule
    메모리 & api 76 $\rightarrow$ 77 MB, web 129 $\rightarrow$ 130 MB \\
    \bottomrule
  \end{tabularx}
\end{table}

\begin{figure}[H]
  \centering
  \begin{tikzpicture}
    \begin{axis}[
      ybar, width=0.62\linewidth, height=5.4cm, bar width=24pt,
      xtick={1,2,3}, xticklabels={전체,web,api}, xmin=0.4, xmax=3.6,
      ymin=0, ymax=110, ymajorgrids=true, grid style={black!12},
      ylabel={p95 latency (ms)}, nodes near coords,
      tick label style={font=\footnotesize}, label style={font=\footnotesize},
      every node near coord/.append style={font=\footnotesize},
      axis lines=left,
    ]
      \addplot[fill=cfcol!60, draw=cfcol] coordinates {(1,82) (2,88) (3,73)};
    \end{axis}
  \end{tikzpicture}
  \caption{L1의 p95 응답 시간. 세 값 모두 100 ms 아래다.
  (src: deploy/traffic/review/02-blocking.md §L1)}\label{fig:07-l1-p95}
\end{figure}

\begin{table}[H]
  \caption{L1 통과 기준과 측정.}\label{tab:07-l1-pass}
  % src: traffic.design.md §10.5 L1; deploy/traffic/review/02-blocking.md §L1
  \small
  \begin{tabularx}{\linewidth}{L{5.2cm} Y L{1.6cm}}
    \toprule
    기준 & 측정 & 판정 \\
    \midrule
    응답 200 $\geq$ 99\,\% & 실패 0 & 통과 \\
    \midrule
    \cfs{HIT} $\geq$ 99\,\% & web \cfs{HIT} 320, api \cfs{HIT} 320 & 통과 \\
    \midrule
    blog-web의 그 path 줄 $\leq$ 10 & 0 (시험 UA 줄 없음) & 통과 \\
    \midrule
    재시작 횟수 불변 & api 5, web 3 (전후 같음) & 통과 \\
    \bottomrule
  \end{tabularx}
\end{table}

\paragraph{읽는 법.}
익명 GET 폭주는 Cloudflare가 모두 받는다.
VM은 아무것도 하지 않았다.

%% ---------------------------------------------------------------------------
\section{결과 2: L2 web, 오리진 \texorpdfstring{\cfs{MISS}}{MISS} 경로}\label{sec:07-result-l2web}

\begin{inbrief}
  \item[Why] 캐시를 놓친 요청을 웹 오리진이 초당 몇 번까지 처리하는지 안다.
  \item[When] run \texttt{38466353765}, 2026/10/13 20:53--20:57 UTC.
  \item[Where] \code{my.prettylog.com/?lt=<run>-<n>}. 매번 \cfs{MISS}다.
  \item[How] 0.5, 1, 2, 4 rps 단계. 중단 조건을 넘으면 멈춘다.
\end{inbrief}

% src: deploy/traffic/review/02-blocking.md §L2: 오리진 MISS 경로 (run 38466353765, 중단)
\paragraph{결과.}
4 rps 단계에서 중단했다.
웹 오리진은 \cfs{MISS} 요청을 초당 2번까지 버티고, 4번에서 시간 초과가 생긴다.
첫 C\textsubscript{web} 값은 약 2 rps다.
재시작이나 5xx는 없었다.
서버는 죽지 않고 느려졌다.

\begin{table}[H]
  \caption{L2 web 단계별 blog-web 로그의 시험 요청 수.}\label{tab:07-l2-stages}
  % src: deploy/traffic/review/02-blocking.md §L2 단계 표
  \small
  \begin{tabularx}{\linewidth}{L{3cm} L{3cm} Y}
    \toprule
    단계 (web) & 시각 (UTC) & blog-web 로그의 시험 요청 \\
    \midrule
    0.5 rps & 20:53 & 13 \\
    \midrule
    1 rps & 20:54 & 44 \\
    \midrule
    2 rps & 20:55 & 87 \\
    \midrule
    4 rps & 20:56 & 172 \\
    \midrule
    중단 & 20:57 & 24 \\
    \bottomrule
  \end{tabularx}
\end{table}

단계별 로그 줄의 합은 340이고 시험 요청 web 320과 다르다.
소스는 이유를 적지 않는다.
세는 범위가 다른 것으로 보이며, 사전 확인 요청이 로그 줄에 들어갔을 수 있다(추정, 확인 안 함).
% src: deploy/traffic/review/02-blocking.md §L2 단계 표 (13+44+87+172+24 = 340), §L2 결과 표 (web 320)

\begin{figure}[H]
  \centering
  \begin{tikzpicture}
    \begin{axis}[
      ybar, width=0.86\linewidth, height=6cm, bar width=22pt,
      xtick={1,2,3,4,5}, xmin=0.4, xmax=5.6,
      xticklabels={{0.5 rps\\20:53},{1 rps\\20:54},{2 rps\\20:55},{4 rps\\20:56},{중단\\20:57}},
      xticklabel style={align=center, font=\footnotesize},
      ymin=0, ymax=200, ymajorgrids=true, grid style={black!12},
      ylabel={Requests in blog-web log (count)}, nodes near coords,
      label style={font=\footnotesize}, tick label style={font=\footnotesize},
      every node near coord/.append style={font=\footnotesize},
      axis lines=left,
    ]
      \addplot[fill=vmcol!55, draw=vmcol] coordinates {(1,13) (2,44) (3,87) (4,172) (5,24)};
    \end{axis}
  \end{tikzpicture}
  \caption{L2 web, 분 단위 blog-web 로그의 시험 요청 수. 단계가 오를수록 늘고, 중단 분에는 24건만 남았다.
  (src: deploy/traffic/review/02-blocking.md §L2)}\label{fig:07-l2-web-chart}
\end{figure}

\begin{table}[H]
  \caption{L2 web 전체 결과.}\label{tab:07-l2-web}
  % src: deploy/traffic/review/02-blocking.md §L2 결과 표
  \small
  \begin{tabularx}{\linewidth}{L{3.2cm} Y}
    \toprule
    항목 & 결과 \\
    \midrule
    요청 & web 320 (api 단계는 시작 전에 중단) \\
    \midrule
    실패 & web 10건 (3.1\,\%), 모두 10 s 시간 초과. 5xx 0 \\
    \midrule
    p95 / 중앙값 & 874 ms / 589 ms \\
    \midrule
    중단 이유 & \code{http_req_failed{target:web}} $>$ 2\,\% (사전 등록 기준) \\
    \midrule
    PM2 재시작 & 없음 (api 5, web 3) \\
    \midrule
    메모리 & web 130 $\rightarrow$ 220 MB (cap 320M), VM available 425 $\rightarrow$ 287 MiB, swap 190 MiB \\
    \bottomrule
  \end{tabularx}
\end{table}

\begin{callout}[새로 드러난 위험]
Cloudflare 한도(처음 가안: IP당 100 / 10 s, 곧 10 rps)는 웹 오리진 용량(약 2 rps)보다 크다.
매번 다른 쿼리(\texttt{/?x=random})로 \cfs{MISS}를 만들면 IP 하나로도 웹 오리진을 밀 수 있다.
이 결과가 한도를 50으로 낮춘 \decision{68}의 이유다.
\end{callout}

%% ---------------------------------------------------------------------------
\section{결과 3: L2-api}\label{sec:07-result-l2api}

\begin{inbrief}
  \item[Why] L2 web이 중단돼 API 단계는 돌지 못했다. C\textsubscript{api}가 없으면 Express 한도의 위쪽 한계를 정할 수 없다.
  \item[When] run \texttt{38468411633}, 2026/10/13 21:09--21:12 UTC. 오너가 승인했다.
  \item[Where] \code{blog-api.prettylog.com/v1/api/posts?limit=1&lt=<run>-<n>}. 매번 \cfs{MISS}다.
  \item[How] 1, 2, 4, 8 rps 단계, 각 30 s. 450건.
\end{inbrief}

% src: deploy/traffic/review/02-blocking.md §L2-api (run 38468411633, 통과, 2026-10-13 21:09-21:12 UTC)
\paragraph{결과.}
통과했다.
450건 모두 \cfs{MISS}로 Express까지 왔다.
실패와 5xx는 0이었다.
API 오리진은 초당 8번에서도 느려지지 않았다.
C\textsubscript{api}는 8 rps 이상이다.
시험 상한이 8 rps라서 실제 한계는 찾지 못했다.

\begin{table}[H]
  \caption{L2-api 결과.}\label{tab:07-l2-api}
  % src: deploy/traffic/review/02-blocking.md §L2-api 표
  \small
  \begin{tabularx}{\linewidth}{L{3.6cm} Y}
    \toprule
    항목 & 결과 \\
    \midrule
    단계 & 1 $\rightarrow$ 2 $\rightarrow$ 4 $\rightarrow$ 8 rps, 각 30 s \\
    \midrule
    요청 & 450, 모두 \cfs{MISS} (Express까지 옴) \\
    \midrule
    실패 / 5xx & 0 / 0 \\
    \midrule
    응답 시간 & 중앙값 501 ms, p95 626 ms, 최대 837 ms \\
    \midrule
    PM2 재시작 & 없음 (api 7, web 5) \\
    \midrule
    메모리 & api 106 MB, web 179 MB, VM available 354 MiB, swap 110 MiB \\
    \midrule
    Express \texttt{rate\_limit} WARN & \textbf{2줄} (21:11:07, 21:11:41), 각 \texttt{bucket=read count=121 limit=120} \\
    \bottomrule
  \end{tabularx}
\end{table}

\begin{figure}[H]
  \centering
  \begin{tikzpicture}
    \begin{axis}[
      ybar, width=0.7\linewidth, height=5.8cm, bar width=20pt,
      xtick={1,2}, xticklabels={중앙값,p95}, xmin=0.4, xmax=2.6,
      ymin=0, ymax=1000, ymajorgrids=true, grid style={black!12},
      ylabel={Latency (ms)}, nodes near coords,
      legend style={at={(0.5,1.04)}, anchor=south, legend columns=-1, draw=none, font=\footnotesize},
      tick label style={font=\footnotesize}, label style={font=\footnotesize},
      every node near coord/.append style={font=\footnotesize},
      axis lines=left,
    ]
      \addplot[fill=vmcol!55, draw=vmcol] coordinates {(1,589) (2,874)};
      \addplot[fill=dbcol!55, draw=dbcol] coordinates {(1,501) (2,626)};
      \legend{L2 web (최대 4 rps), L2-api (최대 8 rps)}
    \end{axis}
  \end{tikzpicture}
  \caption{L2 web과 L2-api의 응답 시간. API는 속도가 두 배인데도 더 빠르다.
  (src: deploy/traffic/review/02-blocking.md §L2 §L2-api)}\label{fig:07-l2-compare}
\end{figure}

\begin{itemize}[nosep]
  \item 웹 오리진(C\textsubscript{web} 약 2 rps)이 먼저 무너진다. 막아야 할 곳은 웹 \cfs{MISS}다.
  \item L-b \texttt{log} 모드가 설계대로 동작했다. 한 창(60 s)에 한 IP가 120번을 넘자 WARN이 창마다 1줄만 나왔고, 요청은 막지 않았다.
\end{itemize}

%% ---------------------------------------------------------------------------
\section{결과 4: Cloudflare 요청 수 제한 발동}\label{sec:07-result-la}

\begin{inbrief}
  \item[Why] 규칙을 켰다는 것과 규칙이 발동한다는 것은 다르다.
  \item[When] 2026/10/13 21:00 UTC(한도 100)와 21:21 UTC(한도 50). 둘 다 Mac에서 \code{curl}로 했다.
  \item[Where] \code{my.prettylog.com}. HIT 경로(\code{/})다.
  \item[How] 동시 요청 25로 짧은 시간에 많이 보내고 429가 나오는지 본다.
\end{inbrief}

% src: deploy/traffic/review/02-blocking.md §L-a 한도 100; §L-a 한도 50 확인 (#68)
\paragraph{결과.}
한도 100에서 250번 중 429 18건이 나왔다.
한도 50에서 150번 중 429 97건이 나왔다.
두 경우 모두 직후에는 200으로 돌아왔다.

\begin{table}[H]
  \caption{Cloudflare ``web rate limit'' 발동 시험.}\label{tab:07-la}
  % src: deploy/traffic/review/02-blocking.md §L-a, §L-a 한도 50 확인
  \small
  \begin{tabularx}{\linewidth}{L{1.8cm} L{5.4cm} L{2.3cm} Y}
    \toprule
    한도 & 시험 & 결과 & 읽는 법 \\
    \midrule
    100 & \code{/robots.txt?rl=<n>} 120번, 동시 20 & 모두 200 & 발동 안 함 \\
    \midrule
    100 & \code{/} (\cfs{HIT}) 250번, 6초, 동시 25 & 200 232, \textbf{429 18} & HIT 요청도 센다. 약 230번째부터 막혔다 \\
    \midrule
    100 & 직후와 12초 뒤 \code{/} & 200 & 차단은 짧게 풀렸다 (10 s) \\
    \midrule
    50 & \code{/} (\cfs{HIT}) 150번, 2초, 동시 25 & 200 \textbf{53}, 429 \textbf{97} & 53번째부터 막혔다 \\
    \midrule
    50 & 15초 뒤 \code{/} & 200 & 차단은 10초 뒤 풀렸다 \\
    \bottomrule
  \end{tabularx}
\end{table}

\begin{figure}[H]
  \centering
  \begin{tikzpicture}
    \begin{axis}[
      ybar, width=0.78\linewidth, height=6.2cm, bar width=24pt,
      xtick={1,2}, xmin=0.4, xmax=2.6, xticklabels={{한도 100\\(250번)},{한도 50\\(150번)}}, xticklabel style={align=center, font=\footnotesize},
      ymin=0, ymax=280, ymajorgrids=true, grid style={black!12},
      ylabel={Responses (count)}, nodes near coords,
      legend style={at={(0.5,1.04)}, anchor=south, legend columns=-1, draw=none, font=\footnotesize},
      tick label style={font=\footnotesize}, label style={font=\footnotesize},
      every node near coord/.append style={font=\footnotesize},
      axis lines=left,
    ]
      \addplot[fill=goodcol!55, draw=goodcol] coordinates {(1,232) (2,53)};
      \addplot[fill=badcol!55, draw=badcol] coordinates {(1,18) (2,97)};
      \legend{200, 429}
    \end{axis}
  \end{tikzpicture}
  \caption{한도 100과 50에서 200과 429의 개수. 한도를 반으로 줄이자 429가 18건에서 97건이 됐다.
  (src: deploy/traffic/review/02-blocking.md §L-a)}\label{fig:07-la-burst}
\end{figure}

\paragraph{읽는 법.}
\begin{itemize}[nosep]
  \item 규칙은 동작한다. 한도 100에서는 세기가 정확하지 않고 늦게 반영됐다. 100번째가 아니라 약 230번째부터 막혔다.
  \item 한도 50에서는 거의 정확했다.
\end{itemize}

\subsection{정상 사용량으로 한도 점검}\label{sec:07-normal-use}

% src: deploy/traffic/review/02-blocking.md §오너 기준 정상 사용량
Cloudflare 한도를 낮추면 정상 사용자가 막히는가?
blog-web 로그(2026/10/13, 시험과 \code{curl} 제외, IP 19개)에서 10초 구간을 세었다.

\begin{table}[H]
  \caption{한도별로 그 한도를 넘었을 10 s 구간 수.}\label{tab:07-normal-use}
  % src: deploy/traffic/review/02-blocking.md §오너 기준 정상 사용량 표
  \small
  \begin{tabularx}{\linewidth}{L{3.2cm} L{4cm} Y}
    \toprule
    한도 (10 s당) & 넘었을 10 s 구간 & 로그인 / 비로그인 \\
    \midrule
    20 & 11 & 7 / 4 \\
    \midrule
    30 & 9 & 5 / 4 \\
    \midrule
    50 & 6 & 3 / 3 \\
    \midrule
    70 & 2 & 2 / 0 \\
    \midrule
    100 & 0 &  \\
    \bottomrule
  \end{tabularx}
\end{table}

\begin{itemize}[nosep]
  \item 최댓값은 10 s에 89번이다. 로그인한 오너의 Mac 브라우저였다.
  \item 원인: Next \code{<Link>}가 화면에 보이는 링크를 미리 받는다(prefetch). 한 번 누를 때 요청이 수십 개 나간다.
  \item 이 표는 VM까지 온 요청만 센다. Cloudflare는 \cfs{HIT}도 세므로 실제 숫자는 더 크다.
  \item 위험: 한도 50에서 정상 사용자가 10 s 동안 막힐 수 있다 (\decision{68}). prefetch 줄이기는 남은 일이다 (\Cref{ch:decisions}).
\end{itemize}

%% ---------------------------------------------------------------------------
\section{결과 5: Express \texttt{log} 모드 확인}\label{sec:07-result-logmode}

\begin{inbrief}
  \item[Why] \texttt{enforce}로 바꾸기 전에 세는 방식이 설계대로인지 본다.
  \item[When] 2026/10/13, 배포 \commit{38b4c12} 뒤와 L2-api run 중이다.
  \item[Where] \texttt{blog-api} 로그의 \texttt{req.ip}, \texttt{req.ssr}, \texttt{rate\_limit} WARN이다.
  \item[How] 로그 필드를 읽고, L2-api에서 일부러 창 한도를 넘긴다.
\end{inbrief}

% src: deploy/traffic/review/02-blocking.md §배포 확인: L-b log 모드, ip, ssr; §L2-api
\paragraph{결과.}
L2-api 중 한 IP가 한 창(60 s)에 121번째 요청을 보내자 WARN 한 줄이 나왔다 (\texttt{bucket=read count=121 limit=120}).
WARN은 창마다 한 줄이고 요청은 막히지 않았다.
기록된 시각은 21:11:07과 21:11:41 두 개다.

\begin{table}[H]
  \caption{배포 확인 (2026/10/13, \texttt{38b4c12}).}\label{tab:07-logmode}
  % src: deploy/traffic/review/02-blocking.md §배포 확인 표
  \small
  \begin{tabularx}{\linewidth}{L{5cm} Y}
    \toprule
    확인 & 결과 \\
    \midrule
    blog-api 로그 \code{jsonPayload.req.ip} & 모든 요청에 있음 \\
    \midrule
    \texttt{ssr=true} & \texttt{profile}, \texttt{tags}, \texttt{posts/users/.../tags} 같은 SSR 호출 \\
    \midrule
    \texttt{ssr=false} & 브라우저와 \code{curl} 요청, 그리고 (당시) SSR의 \texttt{categories} 호출 \\
    \midrule
    \texttt{rate\_limit} WARN & 아직 없음 (L2-api 전) \\
    \midrule
    메모리 (배포 뒤 몇 분) & api 103 MB, web 171 MB, VM available 359 MiB \\
    \bottomrule
  \end{tabularx}
\end{table}

\begin{itemize}[nosep]
  \item 이 시험이 \texttt{ssr=false}로 찍히는 SSR \texttt{categories} 호출을 찾았다. 그 수정이 \Cref{sec:05-clientrequest}(N)이다.
  \item \texttt{enforce}로 바꾸기 전에 모든 SSR 요청이 \texttt{ssr=true}여야 한다. 그렇지 않으면 방문자가 늘 때 사이트가 자기 SSR을 막는다.
\end{itemize}

%% ---------------------------------------------------------------------------
\section{시험 도구에서 찾은 버그}\label{sec:07-k6-bugs}

\begin{inbrief}
  \item[Why] 시험 도구도 코드라서 버그가 있다. 결과를 믿기 전에 도구부터 의심한다.
  \item[When] 첫 L2 시도와 L1 summary를 읽을 때다.
  \item[Where] \code{deploy/traffic/load/l2.js}, \code{common.js}.
  \item[How] 각 버그를 커밋으로 고쳤다. \Cref{tab:07-bugs}가 세 건을 요약한다.
\end{inbrief}

% src: git commits b44b153, 0bf37d0, cbe0f78 (deploy/traffic/load/)
\begin{table}[H]
  \caption{시험 도구의 버그 세 건과 수정.}\label{tab:07-bugs}
  % src: git show b44b153, 0bf37d0, cbe0f78
  \small
  \begin{tabularx}{\linewidth}{L{2.4cm} Y Y}
    \toprule
    커밋 & 버그 & 수정 \\
    \midrule
    \commit{b44b153} & k6가 소수 목표 0.5를 받지 않았다. L2 web의 첫 단계 0.5 rps를 \texttt{target: 0.5}로 적었다 & \texttt{timeUnit}을 \texttt{'2s'}로 바꾸고 목표를 두 배로 적었다 (1, 2, 4, 8). 1은 0.5 rps이고 8은 4 rps다 \\
    \midrule
    \commit{0bf37d0} & Job summary의 요청 수 칸이 \texttt{vals(d).count}를 읽었다 & 요청 수를 \texttt{http\_req\_failed}의 \texttt{passes + fails}로 읽는다 \\
    \midrule
    \commit{cbe0f78} & L2 web이 4 rps에서 중단돼 api 단계가 아예 돌지 못했다 & 세트 \texttt{L2-api}를 추가했다. \texttt{ONLY=api}일 때 web 시나리오와 web 임계값을 지우고 api를 바로 시작한다 \\
    \bottomrule
  \end{tabularx}
\end{table}

수정 전 Job summary 표의 행은 \texttt{| web | undefined | 88 |}처럼 요청 수 칸이 \texttt{undefined}였다(L2 run의 GitHub Actions 로그). \commit{0bf37d0}이 \texttt{common.js}에서 통과와 실패를 함께 세어 고쳤다.
% src: L2 run log (GitHub Actions), commit 0bf37d0

\begin{description}
  \item[0.5 rps] k6의 도착률(arrival-rate) 실행기는 정수 목표를 요구한다. 그래서 ``2초당 1건''으로 쓴다. 속도는 같다. \code{l2.js}의 주석이 이렇게 적는다: \texttt{k6 needs whole-number targets}.
  \item[ONLY=api] 임계값은 데이터가 없는 지표를 평가할 수 없다. web 시나리오를 지웠으면 web 임계값도 함께 지워야 한다.
\end{description}

%% ---------------------------------------------------------------------------
\section{시험으로 알게 된 것}\label{sec:07-conclusions}

\begin{table}[H]
  \caption{시험에서 나온 용량 숫자와 그 쓰임.}\label{tab:07-capacity}
  % src: deploy/traffic/review/02-blocking.md §L2 §L2-api §L-a
  \small
  \begin{tabularx}{\linewidth}{L{3.4cm} L{3.6cm} Y}
    \toprule
    값 & 숫자 & 쓰임 \\
    \midrule
    C\textsubscript{web} & 약 2 rps (첫 값, 판정 아님) & Cloudflare 한도를 정할 때의 위쪽 한계. 웹 \cfs{MISS}를 막아야 한다 \\
    \midrule
    C\textsubscript{api} & 8 rps 이상 (시험 상한) & Express 한도의 위쪽 한계를 정할 때의 근거 \\
    \midrule
    L-a 한도 & 100 $\rightarrow$ 50 / 10 s (\decision{68}) & 웹 오리진 용량보다 큰 한도를 낮춤 \\
    \midrule
    L-b 한도 & 읽기 120 / 60 s, \texttt{log} 모드 & 14일 로그를 본 뒤 확정 \\
    \bottomrule
  \end{tabularx}
\end{table}
